\documentclass[11pt,a4paper]{article}
% Shared preamble for RuPhO English problem statements (match T1 2026 style).
% Use: \input{latex/rupho_en_preamble.tex} after \documentclass.
\usepackage[margin=1in]{geometry}
\usepackage{amsmath,amssymb}
\usepackage{graphicx}
\usepackage{float}
\usepackage{enumitem}
\usepackage{microtype}
\usepackage{caption}
\usepackage{tabularx}
\usepackage{hyperref}

\captionsetup{font=small,labelfont=bf,skip=6pt}

\setlist[itemize]{leftmargin=1.2cm}
\setlist[enumerate]{leftmargin=1.2cm}

% One outer box per subproblem: [id | statement | points] (no inner rules)
\newcommand{\problembox}[3]{%
  \par\addvspace{0.42em}%
  \noindent
  \setlength{\fboxsep}{7.5pt}%
  \setlength{\fboxrule}{0.95pt}%
  \fbox{%
    \begin{minipage}{\dimexpr\linewidth-2\fboxsep-2\fboxrule\relax}
    \begin{tabularx}{\linewidth}{@{}>{\raggedright\arraybackslash\bfseries}p{2.1em} @{\hspace{0.9em}} >{\raggedright\arraybackslash}X @{\hspace{0.9em}} >{\raggedleft\arraybackslash\bfseries}p{2.4em}@{}}
    #1 & #2 & #3
    \end{tabularx}
    \end{minipage}%
  }%
  \par\addvspace{0.32em}%
}


\title{\textbf{Transformers}}
\author{\normalsize X26 (2026) --- English translation}
\date{}
\begin{document}
\maketitle

A transformer is an electromagnetic device with two or more inductively coupled windings. It converts one AC system into another, usually changing voltage and current while keeping frequency constant. In all parts, you may use
\[
\mu_0=4\pi\cdot10^{-7}~\text{H/m}.
\]

Unit vectors along the directions are denoted by the letter c $\wedge$. For example, the unit vector along the $x$ axis is denoted $\hat x$.

Give all numerical answers in SI units.

\section*{Part A. Single-phase transformer (3.6 points)}

The simplest transformer consists of a ferromagnetic core and two windings wound on it with the number of turns $N_1$ and $N_2$. The ohmic resistance of the windings everywhere except part C can be neglected. The magnetic permeability of the core is $\mu\gg 1$, due to which complete magnetic coupling is ensured between the windings of the transformer - that is, we can assume that the magnetic field flux $\Phi$ inside the core is not dissipated.

\section*{Notes}

For the positive direction of magnetic flux in a winding, use the direction generated by the positive current shown in the figure. The mutual inductance of windings $i$ and $j$ is
\[
M=\frac{\partial\Phi_j}{\partial I_i},
\]
where $\Phi_j$ is the flux through winding $j$ and $I_i$ is the current in winding $i$.

\begin{figure}[H]
    \centering
    \includegraphics[width=0.72\textwidth]{T6_EN_assets/fig1.jpg}
    \caption{Figure 1.}
\end{figure}

In points A1 - A3 we find the inductances of the windings $L_1$, $L_2$ and their mutual induction coefficient $M$.

\problembox{A1}{What are the magnetic field fluxes through the primary and secondary windings $\Phi_1$ and $\Phi_2$? Express your answers in terms of $\Phi$, $N_1$ and $N_2$.}{0.20}

To simplify calculations, assume that the cross-sectional area of the core $S$ is constant. The length of the core is $\ell\gg\sqrt{S}$.

\problembox{A2}{Using the circulation theorem for the magnetic field strength vector $\vec H$, find the magnetic field flux $\Phi$ in the core. Express your answer in terms of $S$, $\ell$, $\mu$, $I_1$, $I_2$, $N_1$ and $N_2$.}{0.20}

\problembox{A3}{Obtain expressions for the winding inductances $L_1$, $L_2$ and their mutual inductance coefficient $M$. Express your answer in terms of $S$, $\ell$, $\mu$, $I_1$, $I_2$, $N_1$ and $N_2$.}{0.30}

The interaction of the two windings is completely described by the found quantities $L_{1}, L_{2}$ and $M$. In the diagrams, the transformer is depicted as follows:

\begin{figure}[H]
    \centering
    \includegraphics[width=0.72\textwidth]{T6_EN_assets/fig2.jpg}
    \caption{Figure 2.}
\end{figure}

\problembox{A4}{Find the magnetic induction emf $\mathcal E_{1}$ and $\mathcal E_{2}$ arising in the primary and secondary windings, respectively. Express your answer in terms of $\dot{{I}_{1}}, \dot{{I}_{2}}, L_{1}, L_{2}$ and $M$. Notes: EMF $\mathcal E_{1}$ is considered positive if the electrostatic potential at the top point a1 is greater than at the bottom a2. EMF $\mathcal E_{2}$ is considered positive if the electrostatic potential at the top point b1 is greater than at the bottom b2.}{0.40}

Find the magnetic induction emf $\mathcal E_{1}$ and $\mathcal E_{2}$ arising in the primary and secondary windings, respectively. Express your answer in terms of $\dot{{I}_{1}}, \dot{{I}_{2}}, L_{1}, L_{2}$ and $M$.

A sinusoidal voltage $U_1(t)=\operatorname{Re}\!\left[\tilde U_1 e^{-i\omega t}\right]$ with amplitude $U_0=|\tilde U_1|$ is applied to the primary winding. Use complex amplitudes:
\[
I_1(t)=\operatorname{Re}\!\left[\tilde I_1 e^{-i\omega t}\right],\quad
I_2(t)=\operatorname{Re}\!\left[\tilde I_2 e^{-i\omega t}\right],\quad
U_2(t)=\operatorname{Re}\!\left[\tilde U_2 e^{-i\omega t}\right].
\]

\problembox{A5}{Obtain an expression for the current transformation ratio $k_{I}\equiv{\left|\tilde{I}_{1}\right|}/{\left|\tilde{I}_{2}\right|}$. The response may include $L_1$, $L_2$, $\omega$ and $Z$.}{0.50}

\problembox{A6}{What is the voltage transformation ratio $k_{U}={\left|\tilde{U}_{2}\right|}/{\left|\tilde{U}_{1}\right|}$? The response may include $L_1$, $L_2$, $\omega$ and $Z$.}{0.50}

For power transformers, current is reduced and voltage is increased before transmission to reduce ohmic losses. For such transformers:
\[
k_I,\;k_U\gg 1.
\]

\problembox{A7}{Write down the condition for the number of turns in windings $N_{1}$ and $N_{2}$, under which the expression $(1)$ is satisfied.}{0.30}

Due to the presence of mutual induction, the voltage source $U_1$ can do work despite the absence of active resistance connected directly to it.

\problembox{A8}{Determine the phase difference $\Delta\varphi$ by which the voltage $U_1$ leads the current $I_1$. The answer may include $L_1$, $L_2$, $\omega$ and $Z$.}{0.60}

\problembox{A9}{Determine the time-averaged power $P$ of source $U_1$. The answer may include $U_0$, $Z$, $N_1$ and $N_2$.}{0.60}

\section*{Part B. Instrument current transformer (0.5 points)}

If you've ever done an experiment with electricity, you should know how easily an ammeter can burn out at high currents! To prevent such incidents, a measuring current transformer is used, which reduces the current by a fixed number of times. The primary winding of the transformer is a conductor with the measured alternating current $I$, and an ammeter is connected to the secondary winding with the number of turns $N$. The internal resistance of the ammeter is very small, so in fact the secondary winding of the transformer operates in a mode close to a short circuit. The magnetic permeability of the core is $\mu\gg 1$, its radius is $R$ and is much greater than the thickness of the wire.

\begin{figure}[H]
    \centering
    \includegraphics[width=0.72\textwidth]{T6_EN_assets/fig3.jpg}
    \caption{Figure 3.}
\end{figure}

\problembox{B1}{Find the relation $I / I_{A}$.}{0.50}

\section*{Energy losses in transformers}

To use transformers in practice, losses must be estimated and minimized. Main loss channels are Joule losses in windings, eddy-current losses in the core, and hysteresis losses. In the following parts, treat these contributions as small and independent. Assume equal numbers of turns in primary and secondary windings. For numerical estimates, use
\[
\ell=25~\text{cm},\quad S=100~\text{cm}^2,\quad B_s=1.5~\text{T}.
\]


In the following parts, the influence of each of the factors will be studied, considering these factors to be independent and small, and the efficiency of the transformer will also be calculated. For simplicity, we will assume that the number of turns in the primary and secondary windings is the same.

For numerical calculations, use the following values:

\section*{Part C. Ohmic losses in windings (2 points)}

We will simulate the resistance of the windings by adding a small resistance $r$ to the primary and secondary windings (see figure below).

\begin{figure}[H]
    \centering
    \includegraphics[width=0.72\textwidth]{T6_EN_assets/fig4.jpg}
    \caption{Figure 4.}
\end{figure}

\problembox{C1}{Obtain an expression for the average power $P_0$, which is released at the load if the winding resistance can be neglected. Express your answer in terms of the amplitude of the input voltage $U_0$ and the load resistance $R$. Calculate it for a transformer whose parameters were given in the problem earlier.}{0.20}

\problembox{C2}{Similar to Part A, write equations relating the complex amplitudes $\tilde U_1$, $\tilde U_2$, $\tilde I_1$ and $\tilde I_2$. The response may also include $R$, $r$, single winding inductance $L$ and cyclic frequency $\omega$.}{0.30}

\problembox{C3}{Express $\tilde I_1$ and $\tilde I_2$ in terms of $\tilde U_1$ and other quantities you need, which may include $R$, $r$, $L$ and $\omega$.}{0.40}

\problembox{C4}{Hereinafter, work in the limit $r\ll R, \omega L$. Obtain expressions for the power of Joule losses $P_{r1}$ and $P_{r2}$ released in the primary and secondary windings. Express your answer in terms of the amplitude of the input voltage $U_0$, as well as $r$, $\omega$ and $L$. Calculate the total power of Joule losses on the windings $P_{r} =P_{r1} + P_{r2}$ for the transformer, the parameters of which were given earlier in the problem. Provide detailed calculations for your calculations in the solution sheets.}{1.10}

Obtain expressions for the power of Joule losses $P_{r1}$ and $P_{r2}$ released in the primary and secondary windings. Express your answer in terms of the amplitude of the input voltage $U_0$, as well as $r$, $\omega$ and $L$.

Calculate the total power of Joule losses on the windings $P_{r} =P_{r1} + P_{r2}$ for a transformer whose parameters are given earlier in the problem. Provide detailed calculations for your calculations in the solution sheets.

\section*{Part D. Toki Fuko (2.7 points)}

If an alternating magnetic field is created in a conducting material, an eddy electric field and associated Foucault currents arise in it. This leads to additional energy losses. Let us consider a long cylindrical core with cross-section $S$ and length $\ell$ (we will assume that $\ell\gg\sqrt{S}$). The magnetic field $\vec B(t)=\operatorname{Re}\left[B_0 \hat z e^{-i\omega t}\right]$ penetrates this core along its axis $z$. Neglect the influence of Foucault currents on the magnetic field distribution.

\begin{figure}[H]
    \centering
    \includegraphics[width=0.72\textwidth]{T6_EN_assets/fig5.jpg}
    \caption{Figure 5.}
\end{figure}

\problembox{D1}{Write down the expression for the tangential electric field $E(r,t)$.}{0.30}

\problembox{D2}{Derive an expression for the average power of Joule losses $P_\mathrm{eddy}$ released in the core. Express your answer in terms of the conductivity of the core material $\sigma_\mathrm c$, its cross-sectional area $S$, length $\ell$, and $B_0$ and $f$.}{0.70}

To reduce losses, the core is made layered along its axis - thin layers of non-conducting material are laid between the layers of the ferromagnet. In this case, Foucault currents can only be excited separately in each plate.

Single-phase transformer. Its core is assembled from plates isolated from each other to reduce losses from Foucault currents

\begin{figure}[H]
    \centering
    \includegraphics[width=0.72\textwidth]{T6_EN_assets/fig6.jpg}
    \caption{Figure 6.}
\end{figure}

Let us consider such a plate with thickness $h$ in an external field $\vec B(t)=\operatorname{Re}\left[B_0 \hat y e^{-i\omega t}\right]$ directed parallel to it.

\begin{figure}[H]
    \centering
    \includegraphics[width=0.72\textwidth]{T6_EN_assets/fig7.jpg}
    \caption{Figure 7.}
\end{figure}

In this geometry, the induced electric field is parallel to the plate and perpendicular to the magnetic field. Seek it in the form
\[
\vec E(z,t)=\operatorname{Re}\!\left[A z\,\hat x\,e^{-i\omega t}\right].
\]

\problembox{D3}{Find the factor $A$.}{0.30}

\problembox{D4}{Write down an expression for the power $\mathrm dP/\mathrm dV$ released in such a plate per unit volume, averaged over the entire height of the plate, as well as over time. The response may include $S$, $\ell$, $B_0$, $h$, $\sigma_{\mathrm c}$ and $f$. Let now the core be assembled from plates with the parameters indicated above in the problem. The load $R$ is still connected to the secondary winding. Express the power $P'_\mathrm{eddy}$ that would be released in the core due to Foucault currents in terms of the parameters given in the problem. Calculate this power.}{1.40}

\section*{Part E. Hysteresis (1.2 points)}

Finally, let's look at the effect of hysteresis. We assume that under the influence of the magnetic field strength created by the current in the windings, the magnetic field induction inside the core describes the magnetization curve shown in the figure below. The width of this curve along $H$ is constant, the magnetic field varies in the range from $-B_\mathrm s$ to $+B_\mathrm s$ (the so-called saturation induction) and is zeroed at $\pm H_\mathrm c$ (the so-called coercive force).

\begin{figure}[H]
    \centering
    \includegraphics[width=0.72\textwidth]{T6_EN_assets/fig8.jpg}
    \caption{Figure 8.}
\end{figure}

\problembox{E1}{For simplicity, consider a transformer without load. Record the energy that the source imparts to the core when current $I$ flows through the primary winding and the flux through the core changes to $\mathrm d\Phi$. Then express this energy in terms of $H$, $\mathrm dB$ and core parameters.}{0.40}

\problembox{E2}{We will assume that the last expression is valid in an arbitrary case. Write down the expression for the average power $P_{\mathrm{hyst}}$ hysteresis losses. Express your answer in terms of $H_\mathrm c$, $B_\mathrm s$ and the quantities given at the beginning of part C. Also get the numerical value.}{0.30}

\problembox{E3}{Assuming that all contributions calculated in parts C, D and E are small and independent, calculate the final efficiency of the transformer $\eta$.}{0.50}

\vspace{1em}
\noindent\textit{Source:} \url{https://pho.rs/p/4797}

\end{document}